\chapter{Project 2: Make It Fast, Make It Fit, Make It Scale}\label{ch:proj2}

Project~1 (\cref{ch:proj1}) ended with a transformer language model that trains. This project
turns that model into a \emph{system}: you will measure it until you know where every millisecond
and every byte goes, cut its memory with recomputation and fused kernels, write FlashAttention-2 in
Triton, and then spread training over several GPUs three ways (replicated data parallelism, sharded
optimizer state, and fully sharded parameters), before predicting when each strategy runs out of
steam. The project closes with a speed challenge: the fastest correct training step you can build
for a deliberately oversized configuration.

Like every project in Part~VIII, this chapter is a specification and a guide. It says what to build,
how to set it up in this repository, which chapters teach each piece, what ``working'' looks like,
and how to present the results. It deliberately stops short of solutions: the point is to write the
code yourself. The book's labs implement many of the same ideas against toy settings; use them to
check your understanding, not as a substitute for building your own.

\begin{figure}[t]
\centering
\resizebox{\linewidth}{!}{%
\begin{tikzpicture}[
  part/.style={draw, rounded corners=3pt, minimum height=1.1cm, text width=2.5cm, align=center,
    font=\small, line width=0.6pt},
  tag/.style={circle, fill=sysblue, text=white, font=\bfseries\scriptsize, inner sep=1.5pt},
  arr/.style={-{Stealth[length=2mm]}, line width=0.7pt, draw=black!60}]
\node[part, fill=fdata!15, draw=keygreen] (A) at (0,0) {Measure:\\timing, profiles};
\node[part, fill=fdata!15, draw=keygreen] (B) at (3.2,0) {Precision \&\\memory profiles};
\node[part, fill=fmemory!18, draw=napkin] (C) at (6.4,0) {Residuals \&\\checkpointing};
\node[part, fill=fcompute!15, draw=sysblue] (D) at (9.6,0) {FlashAttention-2\\PyTorch + Triton};
\node[part, fill=fcomm!14, draw=labpurple] (E) at (9.6,-1.9) {Data parallel:\\naive $\rightarrow$ bucketed};
\node[part, fill=fcomm!14, draw=labpurple] (F) at (6.4,-1.9) {Sharded optimizer\\\& fully sharded};
\node[part, fill=black!6, draw=black!60] (G) at (3.2,-1.9) {Cost models:\\DP, FSDP, TP, 2D};
\node[part, fill=warnred!10, draw=warnred] (H) at (0,-1.9) {Speed challenge};
\foreach \n/\k in {A/A,B/B,C/C,D/D,E/E,F/F,G/G,H/H} \node[tag] at (\n.north west) {\k};
\draw[arr] (A)--(B); \draw[arr] (B)--(C); \draw[arr] (C)--(D); \draw[arr] (D)--(E);
\draw[arr] (E)--(F); \draw[arr] (F)--(G); \draw[arr] (G)--(H);
\end{tikzpicture}}
\caption{The eight parts of Project~2. A--D run on one GPU; E--F need two or more (or CPU
processes for correctness); G is analysis, validated against your measurements; H combines
everything.}
\label{fig:proj2:parts}
\end{figure}

% =====================================================================
\section{What you will build}\label{sec:proj2:overview}

\begin{table}[h]
\centering\small
\begin{tabularx}{\linewidth}{@{}lXl@{}}
\toprule
part & deliverable & taught in\\
\midrule
A & A configurable benchmarking script; Nsight Systems profiles with NVTX ranges; a kernel-time
    breakdown of forward, backward and optimizer & \cref{sec:perf:timing,sec:perf:nsys,sec:perf:anatomy}\\
B & Mixed-precision benchmarking; memory snapshots of inference and training; peak-memory
    tables & \cref{sec:perf:amp,sec:perf:memory}\\
C & A residual-accounting tool; a measured checkpointing policy for your model &
    \cref{sec:perf:residuals,sec:perf:ckpt}\\
D & FlashAttention-2 in PyTorch, then Triton forward (with causal masking) and backward;
    attention benchmarks & \cref{sec:attention:flash,sec:triton:flash,sec:perf:fa2}\\
E & Data-parallel training: naive, flattened, overlapped, bucketed; communication benchmarks &
    \cref{sec:dist:comm,sec:dist:dp,sec:sharding:ddp}\\
F & A sharded optimizer wrapper and a fully sharded data-parallel wrapper; memory accounting &
    \cref{sec:dist:zero,sec:sharding:zero,sec:sharding:fsdp}\\
G & Compute-versus-communication models for DP, FSDP, TP and FSDP$\times$TP, checked
    against measurements & \cref{sec:dist:tp,sec:sharding:accounting}\\
H & The fastest correct training step for the challenge configuration &
    \cref{sec:perf:playbook} and all of the above\\
\bottomrule
\end{tabularx}
\caption{Project~2 at a glance.}
\label{tab:proj2:parts}
\end{table}

Learning goals: by the end you should be able to predict a step's time and memory before running
it, explain every large block in a profile, write and test a non-trivial GPU kernel with a custom
backward, implement the communication patterns of data and sharded parallelism from collectives,
and choose a parallel layout from first principles.

% =====================================================================
\section{Setup}\label{sec:proj2:setup}

\subsection{Repository layout}

All project infrastructure lives in \texttt{code/projects/p2\_systems/}:

\begin{lstlisting}[style=plain]
code/projects/p2_systems/
  env_check.py        record versions, GPUs, topology, nsys/ncu, matmul sanity check
  sweep.py            run a command over a JSON grid, one JSON line per run (OOMs too)
  tables.py           JSONL -> Markdown/LaTeX tables for the writeup
  p2/impl.py          YOUR code: the interface the acceptance suite calls (stubs)
  acceptance_p2.py    correctness checks for p2/impl.py (CPU-friendly)
  test_p2_infra.py tests of the infrastructure scripts
\end{lstlisting}

Put your own modules wherever you like (a package next to \texttt{p2/}, or inside your Project~1
package) and make \texttt{p2/impl.py} import from them. The model you profile is your Project~1
transformer; if you skipped that project, use \code{llmsys.gpt.GPTModel} with the sizes below.

\subsection{Environment}

\begin{lstlisting}[style=shell]
cd code
python -m venv .venv && source .venv/bin/activate
pip install -r requirements.txt            # torch, triton (Linux CUDA builds), numpy, ...
python projects/p2_systems/env_check.py            # check the GPU, CUDA, NCCL, nsys and ncu
python -m pytest -q projects/p2_systems            # infrastructure self-test
python -m pytest -q projects/p2_systems/acceptance_p2.py   # fails until you implement p2/impl.py
\end{lstlisting}

Install Nsight Systems and Nsight Compute from NVIDIA's developer site (or your distribution's CUDA
repository) on the GPU machine, and the Nsight Systems GUI on your laptop for viewing reports.
Record \code{env\_check.py --json} with every results file.

\subsection{Model sizes and default settings}

Use these GPT-2-shaped configurations, with a 10{,}000-token vocabulary, batch size 4 and context
length 512 unless a part says otherwise. They are the \code{SIZES} table of Lab~16's
\code{bench.py}, so the two can share code.

\begin{table}[h]
\centering\small
\begin{tabular}{lrrrrr}
\toprule
name & $\dmodel$ & $\dff$ & layers & heads & parameters (approx.)\\
\midrule
small  & 768  & 3072 & 12 & 12 & 0.09\,B + embeddings\\
medium & 1024 & 4096 & 24 & 16 & 0.30\,B + embeddings\\
large  & 1280 & 5120 & 36 & 20 & 0.71\,B + embeddings\\
xl     & 1600 & 6400 & 48 & 25 & 1.47\,B + embeddings\\
\bottomrule
\end{tabular}
\caption{Model sizes for Parts A--F (non-embedding parameters $\approx12L\dmodel^2$).}
\label{tab:proj2:sizes}
\end{table}

\subsection{Hardware and budget}

Parts A--D need one GPU; the bigger sizes and long-context attention benchmarks need 40--80\,GB.
Parts E--F need at least two GPUs on one node for the measurements; their correctness can be
developed entirely on a CPU with the gloo backend, which is how the acceptance suite runs. Part~H is
designed for two 80\,GB GPUs. Before starting, estimate your GPU-hours: list every measurement in
your plan, its configurations and repetitions, and multiply by a predicted step time
(\cref{sec:perf:timing}'s napkin method). Keep the plan in your log and compare at the end; the gap is
itself a useful number.

\subsection{The acceptance interface}

\code{p2/impl.py} defines seven functions; their docstrings are the contract.

\begin{table}[h]
\centering\small
\begin{tabularx}{\linewidth}{@{}lX@{}}
\toprule
function & contract (abridged)\\
\midrule
\ident{flash_attention_pytorch(q, k, v, causal)} & differentiable tiled attention in PyTorch; no $T\times T$ residuals\\
\ident{flash_attention_triton(q, k, v, causal)} & the same through your Triton kernels\\
\ident{make_data_parallel(module)} & wrapper that broadcasts rank 0's weights; \code{sync\_grads()} averages gradients\\
\ident{make_sharded_optimizer(params, cls, **kw)} & optimizer with state split across ranks, same updates as \code{cls}\\
\ident{sharded_state_numel(opt)} & optimizer-state elements stored on this rank\\
\ident{make_fully_sharded(module, compute_dtype)} & parameters stored as shards; gathered just in time\\
\ident{gather_full_state(wrapper)} & full fp32 parameters by name, on every rank\\
\bottomrule
\end{tabularx}
\caption{The Project~2 acceptance interface (\texttt{code/projects/p2\_systems/p2/impl.py}).}
\label{tab:proj2:interface}
\end{table}

The suite checks:
\begin{itemize}
\item attention: forward and all three input gradients against a reference, causal and non-causal,
  and that the bytes saved for backward stay within a budget linear in $T$ (so a saved $T\times T$
  matrix fails);
\item data parallel: two processes with \emph{different} initial weights train for three steps on
  halves of each batch; every rank must end with the parameters of single-process training on the full
  batch;
\item sharded optimizer: the same equivalence for AdamW, plus a bound on the optimizer state stored
  per rank;
\item fully sharded: the same equivalence through \code{gather\_full\_state}, plus a bound on the
  parameter elements each rank stores.
\end{itemize}
Run one part at a time with pytest's \code{-k} option (\code{flash}, \code{data\_parallel},
\code{sharded}, \code{fully}). The environment variable \code{P2\_IMPL} points the suite at another
implementation file.

\begin{figure}[t]
\centering
\resizebox{\linewidth}{!}{%
\begin{tikzpicture}[font=\small,
  box/.style={draw, rounded corners=3pt, align=center, minimum height=0.9cm, text width=3.0cm},
  arr/.style={-{Stealth[length=2mm]}, line width=0.7pt, draw=black!60}]
\node[box, fill=fdata!15, draw=keygreen] (impl) at (0,0) {\texttt{p2/impl.py}\\(your code)};
\node[box, fill=black!6] (suite) at (4.3,0) {\texttt{acceptance\_p2.py}};
\node[box, fill=fcompute!15, draw=sysblue] (ref) at (8.6,1.0) {single-process\\reference training};
\node[box, fill=fcomm!14, draw=labpurple] (ranks) at (8.6,-1.0) {2 gloo processes\\(your wrappers)};
\node[box, fill=fmemory!18, draw=napkin] (cmp) at (12.9,0) {compare parameters\\\& memory bounds};
\draw[arr] (impl)--(suite); \draw[arr] (suite)--(ref); \draw[arr] (suite)--(ranks);
\draw[arr] (ref)--(cmp); \draw[arr] (ranks)--(cmp);
\end{tikzpicture}}
\caption{How the acceptance suite checks distributed code: the same data and optimizer are run
once in a single process on full batches and once across two processes through your wrappers, and
the resulting parameters must agree on every rank.}
\label{fig:proj2:acceptance}
\end{figure}

\subsection{Planning your compute}

Before running anything, write a measurement plan with a predicted cost per entry.
\Cref{tab:proj2:plan} shows the format, filled in with the \emph{method} for each row; put your own
numbers in. Underestimating is normal: profiling runs are slower than timed runs, sweeps include
configurations that run out of memory only after allocating for a while, and every kernel bug costs a
debugging session.

\begin{table}[h]
\centering\small
\begin{tabularx}{\linewidth}{@{}lXl@{}}
\toprule
measurement & how to estimate & GPUs\\
\midrule
A: step-time sweep & configs $\times$ (warm-up + measured steps) $\times$ predicted step time & 1\\
A: Nsight profiles & profiles $\times$ (2--5$\times$ an unprofiled run) + report download time & 1\\
B: memory snapshots & one step per snapshot; viewer time is the real cost & 1\\
C: checkpointing sweep & policies $\times$ (1.33$\times$ step time) $\times$ a few steps & 1\\
D: attention grid & implementations $\times$ $T$ values $\times$ $d$ values $\times$ dtypes $\times$ \code{do\_bench} budget & 1\\
E: collectives & sizes $\times$ world sizes $\times$ backends $\times$ repetitions & 2--8\\
E--F: training variants & variants $\times$ steps $\times$ predicted step time & 2\\
H: optimisation loop & (changes tried) $\times$ (compile + measure); dominated by compile time & 2\\
\bottomrule
\end{tabularx}
\caption{A compute plan template for Project~2.}
\label{tab:proj2:plan}
\end{table}

% =====================================================================
\section{A toolbox for the project}\label{sec:proj2:toolbox}

Three small pieces of infrastructure recur throughout Parts~E--H. They are not part of the
deliverables, but you should have them working before you start Part~E.

\paragraph{Launching processes.} On one node, \code{torchrun} starts one process per GPU and sets
the rank, local rank, world size and rendezvous address in environment variables that
\code{init\_process\_group} reads.
For CPU development, \code{llmsys.dist.launch.run\_distributed} spawns gloo processes from a test
(the acceptance suite uses it). The same script should run in both settings:

\begin{lstlisting}[caption={One entry point for gloo-on-CPU and NCCL-on-GPU runs.},label={lst:proj2:launch}]
import os, torch, torch.distributed as dist

def setup():
    backend = "nccl" if torch.cuda.is_available() else "gloo"
    dist.init_process_group(backend)               # reads RANK/WORLD_SIZE/... from env
    if backend == "nccl":
        torch.cuda.set_device(int(os.environ["LOCAL_RANK"]))
    return dist.get_rank(), dist.get_world_size()

# torchrun --standalone --nproc_per_node=2 train_dp.py --size xl
\end{lstlisting}

\paragraph{Timing distributed code.} Every rank must time the same region, and the step is only as
fast as the slowest rank. Synchronise the device, put a barrier before starting the clock, and
gather all ranks' timings to report the maximum (or the full distribution):

\begin{lstlisting}[caption={Per-rank timing aggregated across ranks.},label={lst:proj2:disttime}]
def timed(fn, iters=10, warmup=5):
    for _ in range(warmup): fn()
    if torch.cuda.is_available(): torch.cuda.synchronize()
    dist.barrier()
    t0 = time.perf_counter()
    for _ in range(iters): fn()
    if torch.cuda.is_available(): torch.cuda.synchronize()   # async_op=False is not enough
    mine = (time.perf_counter() - t0) / iters
    allt = [None] * dist.get_world_size()
    dist.all_gather_object(allt, mine)
    return max(allt), allt
\end{lstlisting}

Note the synchronisation after communication calls: with NCCL, a collective returns once it is
\emph{enqueued} on the GPU stream, even when \code{async\_op=False}.

\paragraph{The cost-model method on another layer.} Part~G asks for compute-versus-communication
models. Here is the method on a single square linear layer $\mat{Y}=\mat{X}\mat{W}$ with
$\mat{X}\in\R^{B\times D}$, $\mat{W}\in\R^{D\times D}$, 2-byte values, in data parallelism over
$N$ devices. \emph{Compute per device:} the backward pass needs $\mat{dX}=\mat{dY}\mat{W}^\top$ and
$\mat{dW}=\mat{X}^\top\mat{dY}$ on $B/N$ rows each, $2\cdot2(B/N)D^2$ FLOPs, taking
$4BD^2/(NC)$ seconds. \emph{Communication:} the gradient $\mat{dW}$ ($2D^2$ bytes) is all-reduced; a
ring all-reduce sends $2\frac{N-1}{N}$ times the data per device (\cref{sec:dist:ring}), taking about
$4D^2/W$ seconds. \emph{Condition:} communication stays hidden while $4D^2/W\le4BD^2/(NC)$, i.e.\
$N\le BW/C$: the global batch in tokens divided by the device's FLOP-to-bandwidth ratio. Larger batches
let you use more devices; the layer width cancels. Part~G repeats this reasoning for the gated MLP,
for sharded parameters (which add all-gathers in forward and backward), and for tensor parallelism
(which communicates activations instead of weights).

% =====================================================================
\section{Part A: Measure before you optimise}\label{sec:proj2:a}

\paragraph{Build.} A benchmarking script that constructs a model from command-line hyperparameters,
creates random token batches, and times forward only, forward and backward, or a full training step
with your optimizer, with a configurable number of warm-up and measured steps. It must synchronise
correctly and report mean and standard deviation. Add NVTX ranges that separate warm-up from
measurement and forward, backward and optimizer from each other, plus an option to annotate the
pieces of attention (scores, softmax, value product).

\paragraph{Measure and explain.}
\begin{enumerate}
\item Forward and backward times for all four sizes. Compare with your napkin predictions and
  explain any factor-of-two discrepancy. How large is the run-to-run variation?
\item Repeat with zero, one and two warm-up steps. Explain each difference in terms of the one-time
  costs listed in \cref{sec:perf:timing}.
\item Profile two sizes at three context lengths of your choice (the largest should be close to the
  memory limit). For each: total forward time versus your timer; the kernel with the most cumulative
  GPU time and how often it runs per forward pass; the non-GEMM kernels that take a visible share
  of time; how the GEMM share changes when the backward pass and optimizer are included.
\item Within attention, compare the time of the softmax with the time of the two matrix products,
  and contrast the ratio with the FLOP ratio.
\end{enumerate}

\paragraph{Sanity checks.} Forward time should scale roughly linearly with tokens and with
non-embedding parameters at short context; backward should be 1.8--2.5$\times$ forward; profiles
captured with an NVTX capture range should contain no compilation or autotuning kernels.

% =====================================================================
\section{Part B: Precision and memory}\label{sec:proj2:b}

\paragraph{Build.} Flags for bf16 autocast and for recording a memory snapshot of the measured
steps (\cref{lst:perf:snapshot}).

\paragraph{Measure and explain.}
\begin{enumerate}
\item A small accumulation experiment of your own design that shows how the precision of a
  running sum (fp32, fp16, bf16, and fp32 accumulation of 16-bit values) affects the result. Relate
  it to how tensor cores and reductions are implemented.
\item For a two-layer toy network with a normalisation layer between the layers, determine
  experimentally (by printing dtypes inside the forward under autocast) the dtype of each parameter,
  activation, the logits, the loss and the gradients. Explain the normalisation layer's treatment and
  whether it would differ between fp16 and bf16.
\item Step times with and without bf16 autocast for every size. Explain the trend with model size.
\item Active-memory timelines for the largest size that fits, for inference only and for a full
  training step, at a short and a long context. Annotate the phases. Tabulate peak memory for each
  case, with and without mixed precision, and explain why mixed precision helps less than ``half''.
\item From the snapshot, identify the largest allocations and the code that makes them. Using NVTX
  ranges in an Nsight memory trace (or saved-tensor hooks), measure how much memory one block keeps
  for backward and how much its gradients occupy, and reconcile both with a calculation.
\end{enumerate}

% =====================================================================
\section{Part C: Residuals and recomputation}\label{sec:proj2:c}

\paragraph{Build.} A residual tracker based on \code{saved\_tensors\_hooks} that counts bytes per
storage and ignores parameters (\cref{lst:perf:tracker} shows the idea; write your own). A way to
checkpoint arbitrary groups of blocks, including nested groups.

\paragraph{Measure and explain.}
\begin{enumerate}
\item The residuals of your normalisation layer and of one full block, eager versus compiled. Explain
  each saved tensor and what fusion removed.
\item Derive the checkpointing policy that minimises peak activation memory when compute is free
  (any nesting allowed): give its asymptotic memory and compute in the number of blocks, with a short
  code sketch. Then, allowing at most one level of recomputation, find the best policy for your
  largest model at a long context, using the cost model of \cref{eq:perf:ckpt} with your own
  measured $a$ and $A$, and confirm it by measuring peak memory for that policy and its neighbours.
\end{enumerate}

\paragraph{Sanity checks.} Checkpointed training must produce gradients identical to
non-checkpointed training (up to dropout, which you should disable for the comparison); one extra
forward per checkpointed block should be visible in the profile.

% =====================================================================
\section{Part D: FlashAttention-2}\label{sec:proj2:d}

\paragraph{Build, in this order.}
\begin{enumerate}
\item A benchmark of your existing attention over a grid of head dimensions and sequence lengths,
  without heads (one head of dimension $d$), timing forward and backward and recording memory before
  the backward. Find where it runs out of memory and account for the memory at that point. Add a
  compiled variant.
\item A \code{torch.autograd.Function} that implements the FlashAttention-2 forward pass in plain
  PyTorch, tile by tile, returning the output and saving $\mat{Q},\mat{K},\mat{V},\mat{O}$ and the
  row log-sum-exp. This is your debugging model.
\item A Triton forward kernel following the same loop structure (one program per query tile and
  batch$\times$head, one loop over key tiles, fp32 running statistics), wrapped in another autograd
  function; then causal masking as an option.
\item A backward pass that recomputes probabilities from the log-sum-exp: first in PyTorch (compiled),
  then optionally in Triton as separate dK/dV and dQ kernels.
\item A benchmark of your kernel against your PyTorch attention with \code{triton.testing.do\_bench}:
  causal, batch 1, sequence lengths from 128 to the largest that fits, head dimensions from 16 to 128,
  bf16 and fp32, reporting forward, backward and end-to-end latency.
\end{enumerate}

\paragraph{Checks.} \code{-k flash} in the acceptance suite. Debug the Triton kernel against your
PyTorch version tile by tile with \code{TRITON\_INTERPRET=1}; compare intermediate values (the running
maximum and sum after each tile) rather than only the final output.

\begin{table}[h]
\centering\small
\begin{tabularx}{\linewidth}{@{}XX@{}}
\toprule
symptom & usual cause\\
\midrule
output correct for the first query tile only & program-id to row-offset arithmetic, or wrong batch stride\\
error grows with $T$ but is small for $T\le$ tile size & running max/sum not rescaled when a new tile raises the max\\
NaNs only with causal masking & a row with no visible key in a tile and $-\infty$ minus $-\infty$; skip or guard\\
correct in fp32, wrong in bf16 & accumulators not fp32, or $P$ not cast before the second dot\\
\code{dV} correct, \code{dQ}/\code{dK} wrong & missing $1/\sqrt d$ scale, or the $\Delta$ term computed from the wrong tensors\\
gradients wrong only for non-contiguous inputs & assumed unit stride; call \code{.contiguous()} or pass all strides\\
\bottomrule
\end{tabularx}
\caption{Kernel debugging: symptoms and their usual causes.}
\label{tab:proj2:kerneldebug}
\end{table}

\begin{pitfall}[Triton versions]
Kernel-language APIs move. The block-pointer API used in many tutorials has been replaced by tensor
descriptors in recent Triton releases (\cref{sec:perf:fa2:blockptr}). Pin the Triton version in your
environment and write explicit pointer arithmetic if you need portability.
\end{pitfall}

% =====================================================================
\section{Part E: Data parallelism}\label{sec:proj2:e}

\paragraph{Build.}
\begin{enumerate}
\item A benchmark of all-reduce latency and bandwidth for tensors from 1\,MB to 1\,GB on 2, 4 and (if
  available) 8 processes, gloo on CPU and NCCL on GPU, following the benchmarking rules of
  \cref{sec:perf:timing} plus warm-up of the communicator and aggregation of timings across ranks.
\item A naive data-parallel trainer: broadcast parameters from rank 0, shard each batch, all-reduce
  each gradient after the backward pass, average, step.
\item A version that flattens all gradients into one buffer and issues a single all-reduce.
\item A wrapper class that launches an asynchronous all-reduce for each parameter as soon as its
  gradient is ready (gradient hooks), and waits for all of them in \code{sync\_grads()} before the
  optimizer step.
\item A bucketed variant with a configurable bucket size.
\end{enumerate}

\paragraph{Measure and explain.} Step time and the fraction spent in communication for each variant
on two GPUs with the xl model; for the bucketed version, a sweep over bucket sizes, compared with a
model of compute and communication time (\cref{sec:dist:dpmodel}). Include profiler screenshots that
show overlap, or its absence.

\paragraph{Checks.} \code{-k data\_parallel}; also verify locally that every rank holds identical
parameters after each step (compare a checksum with an all-gather).

\begin{figure}[t]
\centering
\resizebox{\linewidth}{!}{%
\begin{tikzpicture}[x=0.8cm, y=0.7cm, font=\small]
\node[anchor=east] at (-0.2,2.3) {naive}; \node[anchor=east] at (-0.2,0.8) {overlapped};
\fill[fcompute!45] (0,2) rectangle (5,2.6); \node at (2.5,2.3) {backward};
\fill[fcomm!50] (5,2) rectangle (8.5,2.6); \node at (6.75,2.3) {all-reduce};
\fill[fdata!45] (8.5,2) rectangle (9.5,2.6); \node at (9,2.3) {opt};
\fill[fcompute!45] (0,0.8) rectangle (5,1.4); \node at (2.5,1.1) {backward};
\foreach \x in {1.2,2.4,3.6,4.8} \fill[fcomm!50] (\x,0.1) rectangle ++(1.0,0.5);
\node[font=\scriptsize] at (3.6,-0.25) {bucket all-reduces launched from gradient hooks};
\fill[fdata!45] (5.9,0.8) rectangle (6.9,1.4); \node at (6.4,1.1) {opt};
\draw[keygreen!50!black, thick, -{Stealth}] (9.6,0.5)--(7.0,0.5) node[midway, below, font=\scriptsize] {time saved};
\end{tikzpicture}}
\caption{Why overlap matters (schematic). The naive trainer waits for the whole backward pass before
communicating; the hook-driven version starts reducing each bucket as soon as its gradients exist,
leaving only the last bucket's communication exposed.}
\label{fig:proj2:overlap}
\end{figure}

\begin{table}[h]
\centering\small
\begin{tabularx}{\linewidth}{@{}XX@{}}
\toprule
symptom & usual cause\\
\midrule
hang at the first collective & ranks disagree on world size, address or port; one rank crashed\\
hang after some steps & a collective called by only some ranks (data-dependent branch, unused parameter)\\
parameters drift apart slowly & averaging forgotten or applied twice; different RNG for dropout is fine\\
correct with sync, wrong with async & reading a gradient buffer before \code{wait()}; buffer reused too early\\
works on gloo, fails on NCCL & CPU tensors passed to NCCL, or device not set per rank\\
\bottomrule
\end{tabularx}
\caption{Distributed debugging: symptoms and their usual causes.}
\label{tab:proj2:distdebug}
\end{table}

% =====================================================================
\section{Part F: Sharding}\label{sec:proj2:f}

\paragraph{Build.}
\begin{enumerate}
\item A wrapper that turns any \code{torch.optim.Optimizer} class into a sharded one: each rank owns
  roughly $1/W$ of the parameters' optimizer state, steps only its share, and synchronises updated
  parameters with the other ranks. Support parameter groups and adding groups later.
\item A fully sharded wrapper: each rank stores a shard of every linear and embedding weight (small
  tensors such as norm weights may stay replicated), all-gathers a layer's weights just before its
  forward and backward and frees them afterwards, prefetches the next layer's gather, and
  reduce-scatters gradients into shards. An optional compute dtype casts weights before communication,
  keeping fp32 master shards.
\end{enumerate}

\paragraph{Measure and explain.} Peak memory after model construction, before the optimizer step and
after it, with and without the sharded optimizer, broken down into parameters, gradients, optimizer
state and activations; the time overhead per step; how your sharded optimizer differs from the
partitioning schemes in the literature (\cref{sec:dist:zero}) in memory and communication volume. For
the fully sharded wrapper, the expected and measured memory saving, and a profile showing whether
the weight gathers finish before the layers that need them.

\paragraph{Checks.} \code{-k sharded} and \code{-k fully}.

\begin{figure}[t]
\centering
\resizebox{\linewidth}{!}{%
\begin{tikzpicture}[font=\small,
  st/.style={draw, rounded corners=2pt, minimum height=0.8cm, text width=2.2cm, align=center},
  arr/.style={-{Stealth[length=1.8mm]}, draw=black!60, line width=0.6pt}]
\node[st, fill=fmemory!20] (s0) at (0,0) {shard only\\(fp32, $1/W$)};
\node[st, fill=fcomm!20] (g1) at (2.9,0) {all-gather\\(cast first)};
\node[st, fill=fcompute!20] (f) at (5.8,0) {forward\\of the unit};
\node[st, fill=fmemory!20] (r1) at (8.7,0) {free full\\weights};
\node[st, fill=fcomm!20] (g2) at (8.7,-1.5) {all-gather\\again};
\node[st, fill=fcompute!20] (b) at (5.8,-1.5) {backward\\of the unit};
\node[st, fill=fcomm!20] (rs) at (2.9,-1.5) {reduce-scatter\\gradients};
\node[st, fill=fdata!20] (o) at (0,-1.5) {optimizer step\\on the shard};
\draw[arr] (s0)--(g1); \draw[arr] (g1)--(f); \draw[arr] (f)--(r1); \draw[arr] (r1)--(g2);
\draw[arr] (g2)--(b); \draw[arr] (b)--(rs); \draw[arr] (rs)--(o); \draw[arr] (o)--(s0);
\node[font=\scriptsize, keygreen!50!black] at (4.35,0.75) {prefetch: start next unit's gather here};
\end{tikzpicture}}
\caption{The life of one fully sharded unit during a training step. Only one or two units are
gathered at any time, so peak parameter memory is a shard of everything plus a few full units.}
\label{fig:proj2:fsdp}
\end{figure}

\paragraph{Design decisions to make deliberately.} What is a unit (a block, a linear layer, the
embedding)? How do you pad a flattened unit so that it divides evenly by the world size? How do
you free a gathered tensor without breaking the autograd graph that still refers to it (resizing
the underlying storage to zero and back is the usual trick)? In which order do backward gathers
happen, and how do you know it before the first step? Where exactly does the cast to the compute
dtype happen, and which tensors must never be cast? Write your answers down before coding; they are
the substance of the writeup for this part.

% =====================================================================
\section{Part G: Choosing a parallel layout on paper}\label{sec:proj2:g}

Take the MLP of your model (a gated feed-forward block with weights $\mat{W}_1,\mat{W}_2\in\R^{D\times
D_{\text{ff}}}$ and $\mat{W}_3\in\R^{D_{\text{ff}}\times D}$, batch of $B$ token vectors, 16-bit weights
and activations) on devices with compute rate $C$ FLOP/s and egress bandwidth $W$ bytes/s per device,
in the idealised ring model of \cref{sec:dist:ring}. For each strategy below write the forward and
backward FLOPs per device, the communication time per step, and the condition under which
communication cannot be hidden behind compute. Express the conditions as inequalities on the number
of devices.
\begin{enumerate}
\item Data parallelism over $N_{\text{DP}}$ devices.
\item Fully sharded data parallelism over $N_{\text{FSDP}}$ devices (forward and backward separately).
\item Tensor parallelism over $N_{\text{TP}}$ devices, with the first two weights split by columns and
  the third by rows; write out the sharded backward pass explicitly.
\item The two-dimensional combination $N=N_{\text{FSDP}}N_{\text{TP}}$, first assuming the two axes'
  collectives overlap, then assuming they share the network; find the best split and the largest $N$.
\item Analyse an alternative all-reduce in which each device forwards the \emph{original} chunks
  around the ring and accumulates locally, and compare it with reduce-scatter plus all-gather.
\end{enumerate}
\Cref{sec:sharding:accounting} shows the method on a different layer; reuse its structure. Then
check one prediction against Part~E or F measurements. A model you have tested against
measurements is worth far more than a tidy formula.

% =====================================================================
\section{Part H: The speed challenge}\label{sec:proj2:h}

\paragraph{The task.} Minimise the wall-clock time of one complete training step (forward, loss,
backward, AdamW update) for this configuration on two 80\,GB GPUs:
a pre-norm decoder with RMSNorm, RoPE and a SwiGLU MLP; $\dmodel=2048$, $\dff=6144$, 28 layers,
16 heads, vocabulary 151{,}936, bf16 compute with fp32 master weights, causal attention, two
sequences of 16{,}384 tokens per step (one per GPU).

\paragraph{Rules.} The model's input--output behaviour must not change: logits and loss on a fixed
batch must match a straightforward reference implementation within bf16 tolerance, and so must one
step's parameter update. Everything must be your own code or PyTorch/Triton built-ins. Time with
\code{triton.testing.do\_bench} after warm-up, and report the time from an empty compile/autotune
cache to the first measured step separately (keep it under ten minutes).

\begin{napkinmath}[Does it even fit?]
The configuration is about $1.72\times10^9$ parameters. With fp32 weights, gradients and two Adam
moments that is $16\times1.72\times10^9=27.5$\,GB, or 13.8\,GB per GPU if sharded. Per GPU, one
16{,}384-token sequence through 28 blocks keeps 1.9\,GB of bf16 block inputs if every block is
checkpointed, plus one live block's residuals (about 1.3\,GB at 20 residual units). But the logits
alone are $16{,}384\times151{,}936$ values, 5.0\,GB in bf16 and 10.0\,GB in fp32, before their gradient.
The model fits easily; the loss does not, unless you avoid materialising it (\cref{lst:perf:chunkce}).
The FLOPs, about $6N$ per token plus attention, come to roughly 0.5--0.7\,PFLOP per step depending on
whether causal skipping is counted, i.e.\ a few hundred milliseconds at peak on two H100-class GPUs.
That is the floor to aim at.
\end{napkinmath}

\paragraph{Ideas.} Start from the playbook in \cref{sec:perf:playbook}; profile before and after
every change and keep a table of what each change bought. Likely candidates: tuned FlashAttention
tiles with early exit for causal tiles and separate handling of diagonal tiles; a Triton backward for
attention; a fused or chunked LM head plus cross-entropy; a fused AdamW; selective checkpointing
only where needed; \code{torch.compile} on the blocks; overlapping parameter all-gathers and
gradient reduce-scatters with compute.

\begin{lstlisting}[style=plain,caption={A change log for the speed challenge (one line per experiment).},label={lst:proj2:log}]
id  change                          step_ms  peak_GiB  vs_prev  correct  notes
00  baseline (eager, fp32 loss)     OOM      -         -        -        logits do not fit
01  chunked LM head + loss          ...      ...       -        yes      chunk=1024
02  FA2 kernel, causal early exit   ...      ...       ...      yes
...
\end{lstlisting}

% =====================================================================
\section{Milestones and sanity checks}\label{sec:proj2:milestones}

\begin{table}[h]
\centering\small
\begin{tabularx}{\linewidth}{@{}lX@{}}
\toprule
milestone & you should observe\\
\midrule
A done & timings within $2\times$ of napkin predictions; profiles without warm-up noise\\
B done & autocast speeds up GEMM-heavy sizes most; a sawtooth memory timeline per step\\
C done & identical gradients with checkpointing; measured peak within 10--20\% of the cost model\\
D: PyTorch FA & matches naive attention in fp64 to $10^{-10}$; passes \code{gradcheck}\\
D: Triton FA & acceptance \code{-k flash} passes; memory linear in $T$; beats PyTorch beyond a few thousand tokens\\
E done & \code{-k data\_parallel} passes; overlapped/bucketed variants hide most communication\\
F done & \code{-k sharded}, \code{-k fully} pass; per-rank optimizer state $\approx 1/W$\\
G done & one prediction within a factor of 1.5 of a measurement\\
H done & a correct step, timed, with an itemised list of speedups\\
\bottomrule
\end{tabularx}
\caption{Checkpoints for Project~2.}
\end{table}

\begin{pitfall}[Distributed bugs that tests on two ranks miss]
Uneven shard sizes (parameters whose element count is not divisible by the world size), parameters
that receive no gradient in some steps, gradient accumulation across micro-batches, and
non-deterministic hook order all hide at world size 2 with a toy model. Add your own tests with
world size 3 and with a model that has an unused parameter.
\end{pitfall}

% =====================================================================
\section{Experiment management and the writeup}\label{sec:proj2:writeup}

Generate every table with \code{tables.py} from sweep logs, and every plot with a script stored next to
its data. Name profiles after their configuration (\code{xl\_ctx2048\_bf16\_step.nsys-rep}). For each
part, the writeup should give the question, the setup (hardware, versions, configuration), the
result as a table or figure, and an explanation in terms of the cost models of
\cref{ch:perf,ch:distributed,ch:sharding}. When a result contradicts your model, say so and find out
why.

% =====================================================================
\section{Doing this on less hardware}\label{sec:proj2:small}

\begin{itemize}
\item \textbf{One consumer GPU (16--24\,GB).} Parts A--D with the small and medium sizes and
  contexts up to 2048; the attention benchmarks still show the quadratic wall early. Parts E--F run on
  CPU processes for correctness; measure communication with gloo and treat it as a lower bound.
\item \textbf{Two consumer GPUs without NVLink.} Communication goes over PCIe, so overlap
  matters even more; the cost models of Part~G with a PCIe bandwidth will predict it.
\item \textbf{No GPU.} Everything except Nsight and the performance numbers is possible: the
  acceptance suite runs on CPU with the Triton interpreter and gloo. Use the time to make the
  implementations and the cost models right; rent GPU time for the measurements in one focused
  session using your sweep plan.
\item \textbf{Speed challenge, scaled down.} Halve layers and context for a single 40--48\,GB GPU,
  keep the vocabulary (the loss is the interesting part), and report the scaling you expect for the
  full configuration.
\end{itemize}
